%======================================================================
%  Capitulo 1 - Conocimiento general de la organizacion
%======================================================================

\section{Conocimiento general de la organización}

%-----------------------------------------------------------------------
\subsection{Constitución de la Empresa: Nombre, actividad económica, antecedentes.}

Enlace Express S.A.S. es una empresa de mensajería y paquetería urbana
que se planea constituir en Bogotá D.C. como sociedad por acciones
simplificada. Su servicio consiste en recoger paquetes y documentos en
las instalaciones del cliente, llevarlos al centro de distribución,
clasificarlos por zona de destino y entregarlos al destinatario. Ofrece
tres modalidades según el tiempo de entrega (mismo día, día siguiente y
estándar de dos días hábiles), y cobra según el peso, la zona de destino
y la modalidad.

Sus clientes serán tiendas en línea, droguerías y oficinas que envían
documentos. Para mantener una operación acotada en los primeros años, no
atenderá público ocasional en mostrador ni hará envíos fuera de la
ciudad. Iniciará con un centro de distribución y una flota reducida de
motocicletas, furgonetas y bicicletas eléctricas, y planea abrir más
centros, por lo que su sistema de información debe soportar desde el
comienzo varios centros.

\textbf{Antecedentes del sector.} El crecimiento de las compras por
internet en Colombia multiplicó las empresas pequeñas y medianas de
mensajería urbana. La mayoría nació de forma improvisada y nunca ajustó
su manejo de la información, lo que produce cinco problemas recurrentes:

\begin{enumerate}[label=\alph*),leftmargin=1.2cm,itemsep=0pt]
  \item Cada sede lleva su propio archivo de Excel, así que no se sabe
        cuánto envía en total un cliente atendido desde dos sedes.
  \item No se sabe en qué estado va un envío: hay que llamar al
        mensajero.
  \item La facturación se arma juntando a mano los archivos de cada
        sede, lo que tarda días y genera errores que los clientes
        reclaman.
  \item Los documentos de los vehículos y las licencias se controlan
        aparte de los despachos, y es fácil que salga a ruta un vehículo
        con documentos vencidos.
  \item No hay indicadores confiables de cumplimiento, y se decide por
        percepción.
\end{enumerate}

Estos problemas vienen de crecer sin un sistema único, y corregirlos
después es costoso porque obliga a migrar datos incompletos y a cambiar
hábitos de trabajo. Por eso Enlace Express S.A.S. arrancará desde el
primer día con una base de datos centralizada que los resuelva; ese es
el propósito de este proyecto.

%-----------------------------------------------------------------------
\subsection{Misión de la Empresa}

Prestar un servicio de mensajería y paquetería urbana rápido y confiable
a los comercios y empresas de Bogotá, cumpliendo los tiempos de entrega
acordados y cuidando la mercancía que se nos confía, con personal
capacitado, una flota eficiente y un manejo ordenado de la información.

%-----------------------------------------------------------------------
\subsection{Visión de la Empresa}

Para el año 2032, Enlace Express S.A.S. será una empresa de mensajería
urbana reconocida en Bogotá por su cumplimiento y por la trazabilidad de
sus envíos, con varios centros de distribución, una flota
mayoritariamente eléctrica y un sistema propio que permita a los
clientes consultar el estado de sus envíos en cualquier momento.

%-----------------------------------------------------------------------
\subsection{Objetivos de la Empresa}

\textbf{Objetivo general.} Prestar un servicio de mensajería urbana que
cumpla los tiempos de entrega pactados con cada cliente, garantizando la
integridad de la mercancía y el seguimiento completo del envío desde su
recolección hasta su entrega.

\textbf{Objetivos específicos.}

\begin{enumerate}[label=\textbf{OE\arabic*.},leftmargin=1.4cm,itemsep=0pt]
  \item Cumplir los tiempos de entrega ofrecidos y mantener baja la
        cantidad de envíos con novedades.
  \item Informar oportunamente a los clientes el estado de sus envíos y
        responder rápido sus reclamaciones.
  \item Construir una base de clientes estable, con tarifas ajustadas a
        su volumen y a las zonas que más usan.
  \item Operar desde el primer día sobre un único sistema de
        información, capaz de soportar nuevos centros sin reprocesos.
  \item Mantener un equipo de mensajeros capacitado, vinculado
        formalmente y con poca rotación.
  \item Incorporar progresivamente vehículos eléctricos a la flota.
\end{enumerate}

%-----------------------------------------------------------------------
\subsection{Metas de la Empresa}

Como la empresa aún no opera, las metas se comparan con los valores
habituales en empresas del sector de tamaño similar
(tabla~\ref{tab:metas}).

{\singlespacing
\begin{table}[H]
\centering
\renewcommand{\arraystretch}{1.3}
\begin{tabularx}{\textwidth}{L{1.2cm} X L{2.4cm} L{2.4cm}}
\toprule
\textbf{Obj.} & \textbf{Meta} & \textbf{Referencia del sector} & \textbf{Meta de Enlace Express} \\
\midrule
OE1 & Entregar a tiempo los envíos & 85\% & 95\% \\
OE1 & Mantener bajos los envíos con novedad & 9\% & 5\% \\
OE2 & Responder las reclamaciones & 24 horas & 4 horas \\
OE2 & Ofrecer consulta del estado del envío & Poco frecuente & Desde el primer día \\
OE3 & Alcanzar una base de clientes estable & --- & 120 clientes \\
OE3 & Mantener bajas las facturas con reclamo & 6\% & 1\% \\
OE4 & Operar sobre un sistema único & Archivos por sede & 1 sistema \\
OE5 & Mantener baja la rotación de mensajeros & 38\% & 20\% \\
OE5 & Evitar despachos con documentos vencidos & Sin control & 0 casos \\
OE6 & Incorporar vehículos eléctricos & 10\% & 35\% \\
\bottomrule
\end{tabularx}
\caption{Metas de la empresa frente a los valores de referencia del sector.}
\label{tab:metas}
\end{table}
\par}

%-----------------------------------------------------------------------
\subsection{Políticas de la Empresa}

\begin{itemize}[leftmargin=1.2cm,itemsep=2pt]
  \item \textbf{Calidad del servicio.} Todo envío se registra con los
        datos completos del remitente, el destinatario y el paquete
        antes de despacharse, y ninguno se marca como entregado sin
        prueba de entrega (firma, número de documento o fotografía del
        destinatario).
  \item \textbf{Seguridad vial.} Ningún vehículo sale a un despacho con
        un documento obligatorio vencido, y ningún mensajero sale a ruta
        con la licencia vencida o sin la capacitación anual en seguridad
        vial.
  \item \textbf{Manejo de la mercancía.} El responsable del envío es el
        empleado que lo tiene asignado, y todo cambio de responsable
        queda registrado. Los envíos frágiles o de alto valor solo se
        entregan al destinatario.
  \item \textbf{Datos personales.} Los datos de remitentes,
        destinatarios y empleados se tratan según la Ley 1581 de 2012 y
        solo para prestar el servicio.
  \item \textbf{Acceso a la información.} Cada empleado accede solo a la
        información que necesita, y las creaciones, modificaciones y
        anulaciones importantes quedan registradas.
  \item \textbf{Tarifas.} Todo servicio está amparado por un contrato
        vigente con tarifas por zona, peso y modalidad, y se cobra con
        la tarifa vigente en la fecha del envío.
\end{itemize}

%-----------------------------------------------------------------------
\subsection{Estrategias de la Empresa}

\begin{enumerate}[label=\arabic*.,leftmargin=1.2cm,itemsep=0pt]
  \item \textbf{Zonificación:} dividir Bogotá en zonas asignadas a un
        centro, con mensajeros fijos por zona.
  \item \textbf{Seguimiento por eventos:} registrar cada cambio de
        estado del envío con su fecha y responsable.
  \item \textbf{Consulta en línea:} permitir al cliente consultar su
        envío con el número de guía.
  \item \textbf{Tarifas por perfil:} fijar tarifas según el volumen y
        las zonas de cada cliente.
  \item \textbf{Facturación automática:} facturar directamente desde
        los envíos registrados.
  \item \textbf{Sistema único desde el inicio:} una base de datos
        centralizada, con reglas que impidan registros incompletos y que
        soporte nuevos centros.
  \item \textbf{Control de vencimientos:} validar los documentos del
        vehículo y del mensajero al armar el despacho.
  \item \textbf{Renovación de la flota:} incorporar vehículos eléctricos
        en las zonas de recorridos cortos.
\end{enumerate}

%-----------------------------------------------------------------------
\subsection{Estructura Organizacional de la Empresa (Organigrama)}

La estructura es funcional y tiene tres niveles: la Gerencia General,
cuatro áreas y sus unidades operativas
(figura~\ref{fig:organigrama}). En el primer año varias unidades estarán
a cargo de una sola persona, pero se definen desde ahora porque son
funciones que el sistema debe diferenciar.

{\singlespacing
\begin{figure}[H]
\centering
\begin{tikzpicture}[node distance=0pt]

% --- Nivel 3: unidades operativas ---
\node[caja,text width=3.6cm,minimum height=1.2cm] (cdf)  at (10.6,  0.00) {Centros de\\Distribución};
\node[caja,text width=3.6cm,minimum height=1.2cm] (mesa) at (10.6, -1.45) {Mesa de Servicio};
\node[caja,text width=3.6cm,minimum height=1.2cm] (men)  at (10.6, -2.90) {Mensajeros};

\node[caja,text width=3.6cm,minimum height=1.2cm] (tal)  at (10.6, -4.60) {Taller y\\Mantenimiento};
\node[caja,text width=3.6cm,minimum height=1.2cm] (doc)  at (10.6, -6.05) {Control Documental};

\node[caja,text width=3.6cm,minimum height=1.2cm] (ase)  at (10.6, -7.75) {Asesores Comerciales};
\node[caja,text width=3.6cm,minimum height=1.2cm] (fac)  at (10.6, -9.20) {Facturación y Cartera};

\node[caja,text width=3.6cm,minimum height=1.2cm] (sel)  at (10.6,-10.90) {Contratación};
\node[caja,text width=3.6cm,minimum height=1.2cm] (sst)  at (10.6,-12.35) {Seguridad y Salud\\en el Trabajo};

% --- Nivel 2: areas ---
\node[caja,text width=3.2cm,minimum height=1.2cm] (op)  at (4.8, -1.45) {Área de\\Operaciones};
\node[caja,text width=3.2cm,minimum height=1.2cm] (flo) at (4.8, -5.32) {Área de Flota y\\Mantenimiento};
\node[caja,text width=3.2cm,minimum height=1.2cm] (com) at (4.8, -8.47) {Área Comercial\\y Financiera};
\node[caja,text width=3.2cm,minimum height=1.2cm] (th)  at (4.8,-11.62) {Área de Talento\\Humano};

% --- Nivel 1: gerencia ---
\node[cajaf,text width=3.2cm,minimum height=1.2cm] (ger) at (0, -6.53) {Gerencia General};

% --- Conectores gerencia -> areas ---
\draw[linea] (ger.east) -- (2.6,-6.53);
\foreach \a in {op,flo,com,th}{ \draw[linea] (2.6,-6.53) |- (\a.west); }

% --- Conectores areas -> unidades ---
\foreach \a/\hijos in {op/{cdf,mesa,men}, flo/{tal,doc},
                       com/{ase,fac}, th/{sel,sst}}{
  \draw[linea] (\a.east) -- ++(0.8,0);
  \foreach \h in \hijos { \draw[linea] ($(\a.east)+(0.8,0)$) |- (\h.west); }
}

\end{tikzpicture}
\caption{Organigrama previsto para Enlace Express S.A.S.}
\label{fig:organigrama}
\end{figure}
\par}

{\singlespacing\small
\begin{longtable}{L{4.2cm} L{10.2cm}}
\caption{Responsabilidades de las unidades de la organización.}
\label{tab:responsabilidades}\\
\toprule
\textbf{Unidad} & \textbf{Responsabilidad} \\
\midrule
\endfirsthead
\toprule
\textbf{Unidad} & \textbf{Responsabilidad} \\
\midrule
\endhead
\bottomrule
\endfoot
Gerencia General & Dirige y representa la empresa; aprueba políticas y presupuesto. \\
Área de Operaciones & Planea y controla la recolección, clasificación, despacho y entrega. \\
\quad Centros de Distribución & Reciben, clasifican, arman los despachos y cierran las rutas. \\
\quad Mesa de Servicio & Atiende consultas, registra novedades y gestiona reintentos y devoluciones. \\
\quad Mensajeros & Recogen y entregan los envíos y registran las pruebas de entrega. \\
Área de Flota y Mantenimiento & Mantiene los vehículos disponibles, programa mantenimientos y controla vencimientos de documentos. \\
Área Comercial y Financiera & Consigue y retiene clientes, define tarifas y maneja facturación y cartera. \\
Área de Talento Humano & Selección, contratación, capacitación y seguridad en el trabajo. \\
\end{longtable}
\par}

%-----------------------------------------------------------------------
\subsection{Estructura funcional (Enfoque sistémico o Mapa de procesos).}

Vista como sistema, la empresa transforma solicitudes de envío en
entregas efectivas. Sus \textbf{entradas} son las solicitudes de
recolección, la mercancía, los contratos y tarifas, y los vehículos y el
personal disponibles; sus \textbf{procesos}, el registro, la
clasificación, el enrutamiento, el despacho, la entrega y la
facturación; sus \textbf{salidas}, los envíos entregados con prueba, las
devoluciones, las facturas y los indicadores; y su
\textbf{retroalimentación}, las novedades, reclamaciones e indicadores
que sirven para ajustar zonas y recursos. Los procesos se agrupan en
estratégicos, misionales y de apoyo (figura~\ref{fig:mapaprocesos}).

{\singlespacing
\begin{figure}[H]
\centering
\begin{tikzpicture}

% --- Marcos de las bandas ---
\begin{scope}[on background layer]
  \draw[dashed] (0, 7.6) rectangle (15.8, 11.0);
  \draw[dashed] (0, 3.8) rectangle (15.8,  7.2);
  \draw[dashed] (0, 0.0) rectangle (15.8,  3.4);
\end{scope}

% --- Etiquetas de banda ---
\node[etiqueta] at (0.6, 9.3) {PROCESOS\\ESTRATÉGICOS};
\node[etiqueta] at (0.6, 5.5) {PROCESOS\\MISIONALES};
\node[etiqueta] at (0.6, 1.7) {PROCESOS\\DE APOYO};

% --- Procesos estrategicos ---
\node[caja,text width=3.4cm,minimum height=1.5cm] at ( 3.15,9.3) {Direccionamiento\\estratégico};
\node[caja,text width=3.4cm,minimum height=1.5cm] at ( 8.45,9.3) {Gestión de calidad\\y mejora continua};
\node[caja,text width=3.4cm,minimum height=1.5cm] at (13.75,9.3) {Gestión de\\indicadores};

% --- Procesos misionales ---
\node[caja,text width=2.8cm,minimum height=1.5cm] (m1) at ( 2.85,5.5) {Gestión\\comercial};
\node[caja,text width=2.8cm,minimum height=1.5cm] (m2) at ( 6.58,5.5) {Recolección y\\clasificación};
\node[caja,text width=2.8cm,minimum height=1.5cm] (m3) at (10.32,5.5) {Distribución\\y entrega};
\node[caja,text width=2.8cm,minimum height=1.5cm] (m4) at (14.05,5.5) {Novedades y\\devoluciones};

\foreach \a/\b in {m1/m2, m2/m3, m3/m4}{ \draw[flecha] (\a) -- (\b); }

% --- Procesos de apoyo ---
\node[caja,text width=2.8cm,minimum height=1.5cm] at ( 2.85,1.7) {Gestión de\\flota};
\node[caja,text width=2.8cm,minimum height=1.5cm] at ( 6.58,1.7) {Talento\\humano};
\node[caja,text width=2.8cm,minimum height=1.5cm] at (10.32,1.7) {Tecnología};
\node[caja,text width=2.8cm,minimum height=1.5cm] at (14.05,1.7) {Facturación\\y cartera};

\end{tikzpicture}
\caption{Mapa de procesos previsto para Enlace Express S.A.S.}
\label{fig:mapaprocesos}
\end{figure}
\par}

Los procesos misionales de operación (recolección y clasificación,
distribución y entrega, y novedades y devoluciones) son el centro del
negocio y los que más datos generan; flota y talento humano aportan la
información para validar los despachos, y la facturación usa el
resultado de la operación. Esta relación define las áreas funcionales
del capítulo 2 y es la base para elegir el módulo que se prioriza en la
sección~\ref{sec:modulo}.
